% !TEX program = xelatex
\documentclass[UTF8,a4paper,11pt,fontset=windows]{ctexart}
\usepackage[margin=22mm,headheight=15pt]{geometry}
\usepackage{amsmath,amssymb,graphicx,booktabs,tabularx,array,float,xcolor,listings,fancyhdr}
\usepackage[unicode,colorlinks=true,linkcolor=blue!45!black,urlcolor=blue!45!black]{hyperref}
\usepackage{xurl}
\usepackage{caption}
\usepackage{enumitem}
\setlist{itemsep=2pt,topsep=4pt,parsep=0pt}
\definecolor{accent}{HTML}{174A71}
\definecolor{codebg}{HTML}{F4F6F8}
\setlength{\parindent}{2em}
\setlength{\parskip}{3pt}
\renewcommand{\arraystretch}{1.25}
\captionsetup{font=small,labelfont=bf,skip=5pt}
\ctexset{section={format=\Large\bfseries\color{accent}},subsection={format=\large\bfseries}}
\pagestyle{fancy}
\fancyhf{}
\fancyhead[L]{区块链基础及应用\quad Ex1}
\fancyhead[R]{邹俊轩\quad 2412592}
\fancyfoot[C]{\thepage}
\setmonofont{Consolas}
\lstset{language=Python,basicstyle=\small\ttfamily,keywordstyle=\color{accent}\bfseries,
 commentstyle=\color{gray},stringstyle=\color{teal!65!black},backgroundcolor=\color{codebg},
 frame=single,rulecolor=\color{gray!35},breaklines=true,columns=fullflexible,
 keepspaces=true,showstringspaces=false,tabsize=4,xleftmargin=5pt,xrightmargin=5pt,
 aboveskip=7pt,belowskip=7pt}
\newcommand{\txhash}[1]{\par\noindent\begingroup\small\ttfamily\nolinkurl{#1}\par\endgroup}
\newcommand{\shot}[3]{\begin{figure}[H]\centering\includegraphics[width=\linewidth,height=#2,keepaspectratio]{assets/#1}\caption{#3}\end{figure}}
\hypersetup{pdftitle={区块链基础及应用 Ex1 实验报告},pdfauthor={邹俊轩（2412592）}}
\begin{document}
\begin{center}
{\LARGE\bfseries《区块链基础及应用》实验报告\par}
\vspace{5mm}
{\Large Ex1：比特币测试网 P2PKH 交易\par}
\vspace{5mm}
姓名：\textbf{邹俊轩}\qquad 学号：\textbf{2412592}\par
实验日期：2026 年 9 月 28 日
\end{center}
\vspace{3mm}
\section{项目构成与实验目标}
项目 GitHub 仓库：\url{https://github.com/pointzu/blockchain_experiment_ex1}。源码、交易记录、截图及本报告均整理于该仓库。\par
本实验依据课程 \texttt{Ex1.pdf} 和提交要求，使用 Python 构造 Bitcoin testnet3 交易。在领取测试币后，将一个未花费输出拆分为十个输出，再通过显式编写的 P2PKH 锁定脚本与解锁脚本，花费其中一份并发送至课程配置提供的地址。代码说明、金额和结果均对应本次实际实验。
\subsection{项目文件}
\par\noindent\begin{tabularx}{\linewidth}{@{}>{\ttfamily\small}l X@{}}
\toprule 文件 & 本次用途与修改 \\
\midrule
config.py & 从本地文件或环境变量读取私钥，派生公钥和地址并核对领取记录。\\
keygen.py & 预备实验生成测试网私钥和地址；\\
split\_test\_coins .py & 设置领取交易、输出编号、总输出金额及份数，构造十个独立输出并签名。\\
ex1.py & 补全两个 P2PKH 函数，花费分币交易的第 0 号输出。\\
utils.py & 金额转整数 satoshi、本地脚本验证、链上输入检查、广播和哈希保存。\\
runtime\_setup.py & 为当前 Python 进程优先使用已有的同位数 OpenSSL DLL。\\
transactions.py & 保存领取、分币和发币的真实交易哈希。\\
check\_results.py & 通过公开 API 查询交易确认状态、输入与输出。\\
requirements.txt & python-bitcoinlib 与 requests 依赖。\\
\bottomrule
\end{tabularx}
\subsection{提交材料与完成标准}
提交要求包括填好的脚本源代码、交易输出结果和相关截图。对应保存 \texttt{split\_output.txt}、\texttt{ex1\_output.txt} 和领取、分币、发币截图；\texttt{chain\_status.txt} 为最终确认状态的补充记录。以脚本验证通过、广播成功且链上交易已确认作为完成标准。

\newpage
\section{实验原理}
\subsection{密钥、地址与 UTXO}
私钥用于签名，公钥用于验证签名。本实验使用的传统 P2PKH 地址由公钥的 Hash160 值结合测试网版本信息和校验码编码得到。数字签名提供花费授权；链上数据和收款地址则可以公开查询。

比特币通过未花费交易输出（UTXO）表示可花费资金。一个输入用“上一笔交易的哈希 + 输出编号”定位旧输出。编号从 0 开始，与用户做过多少次交易无关。旧输出被整体花费，新交易创建新的输出；同一个旧输出不能被重复花费。
\subsection{P2PKH 脚本执行}
锁定脚本 \texttt{scriptPubKey} 指定花费条件，解锁脚本 \texttt{scriptSig} 提供签名和公钥：
\begin{lstlisting}[language={}]
scriptSig:    <signature> <public_key>
scriptPubKey: OP_DUP OP_HASH160 <public_key_hash>
              OP_EQUALVERIFY OP_CHECKSIG
\end{lstlisting}
节点在验证旧输出的花费时执行两部分脚本。下表右侧为栈顶：
\par\noindent\begin{tabularx}{\linewidth}{@{}l X X@{}}
\toprule 步骤 & 栈内容 & 作用 \\
\midrule
解锁材料入栈 & 签名、公钥 & 提交签名和公钥。\\
\texttt{OP\_DUP} & 签名、公钥、公钥 & 复制公钥，保留一份用于验签。\\
\texttt{OP\_HASH160} & 签名、公钥、计算出的哈希 & 对公钥执行 SHA-256 后再执行 RIPEMD-160。\\
压入预期哈希 & 签名、公钥、计算出的哈希、预期哈希 & 提供旧输出中指定的公钥哈希。\\
\texttt{OP\_EQUALVERIFY} & 签名、公钥 & 比较哈希；相同则移除二者并继续，不同则失败。\\
\texttt{OP\_CHECKSIG} & 验证成功时为 True & 用公钥验证本笔交易的签名。\\
\bottomrule
\end{tabularx}
\subsection{金额与手续费}
交易输出金额以整数 satoshi（sat）表示，$1\ \mathrm{BTC}=10^8\ \mathrm{sat}$。本实验的金额关系为：
\[
\text{输入总额}=\text{输出总额}+\text{手续费}.
\]
脚本变量 \texttt{amount\_to\_send} 表示输出金额，并非手续费。若未创建找零输出，输入与输出的差额不会自动退回，而是全部成为手续费。本实验先分币再花其中一份，以保留其他九份供后续实验使用。

\newpage
\section{实验步骤与实现}
\subsection{环境配置与领取测试币}
实验在 Windows 和 Python 3.13 环境中运行，依赖 \texttt{python-bitcoinlib} 与 \texttt{requests}。曾因 Python 加载其他软件目录中的不匹配 DLL 出现 \texttt{WinError 193}；通过 \texttt{runtime\_setup.py} 在当前进程的搜索路径中优先放置 \texttt{ObtainTestcoin} 目录，并核对 DLL 位数后，钱包模块正常导入。

将预备实验保留的原私钥写入本地配置后，派生地址与领取截图中的地址一致：
\txhash{mzT3NxNH1iGLvKiHDzJyN4NtLogEFLv3A2}
\shot{faucet.png}{77mm}{水龙头领取成功记录：0.00150524 BTC}
领取交易哈希为：
\txhash{0d9b2e7b232bc18f54d3c86a77cf1900c3cf544f83c0610b083fc69490345560}
公开 API 的查询结果表明，该交易的第 \textbf{1} 号输出属于本人地址，金额为 \textbf{150524 sat}，确认区块为 \textbf{5148688}。因此分币使用 \texttt{utxo\_index=1}。

课程原示例使用 BlockCypher。实验期间该接口查不到领取交易；根据助教提示，使用 testnet3 的 mempool 与 Blockstream 查询。广播实现另行改为 Blockstream 的 Esplora \texttt{POST /tx} 接口，请求正文直接包含十六进制原始交易，成功返回交易哈希。查询网站替换和广播接口改动是两项不同的调整。

\newpage
\subsection{构造并签名分币交易}
为后续练习保留多个独立输出，设置 \texttt{n=10}。分币输入使用领取交易的第 1 号输出，所有新输出仍锁定到本人地址。关键参数为：
\begin{lstlisting}
amount_to_send = Decimal('0.00149')
txid_to_spend = (
    '0d9b2e7b232bc18f54d3c86a77cf1900c3cf544f83c0610b083fc69490345560')
utxo_index = 1
n = 10
\end{lstlisting}
分币交易的关键构造与签名代码如下（摘录自完成的脚本）：
\begin{lstlisting}
txin_scriptPubKey = my_address.to_scriptPubKey()
txin = create_txin(txid_to_spend, utxo_index)
each_amount = Decimal(str(amount_to_send)) / n
tx = CMutableTransaction(
    [txin],
    [create_txout(each_amount, txin_scriptPubKey)
     for _ in range(n)])
sighash = SignatureHash(txin_scriptPubKey, tx, 0, SIGHASH_ALL)
txin.scriptSig = CScript([
    my_private_key.sign(sighash) + bytes([SIGHASH_ALL]),
    my_public_key])
VerifyScript(txin.scriptSig, txin_scriptPubKey,
             tx, 0, (SCRIPT_VERIFY_P2SH,))
\end{lstlisting}
这里针对整笔十输出交易计算签名摘要。\texttt{SIGHASH\_ALL} 承诺交易的所有输出，使签名与本次付款内容关联。代码为每一份创建独立的 \texttt{CMutableTxOut} 对象。

\texttt{utils.py} 中采用 \texttt{Decimal} 计算 satoshi，检查金额没有不足 1 sat 的部分后转为整数，避免浮点误差。链上预检查还验证输入交易已确认、输出属于本地地址且尚未花费。
\[
150524 = 10\times14900+1524\quad (\mathrm{sat}).
\]
本次实际交易大小为 \textbf{497 bytes}，手续费率为 $1524/497\approx3.07\ \mathrm{sat/vB}$。这些数值来源于实际广播输出。

\newpage
\subsection{分币运行结果与链上确认}
先以默认预览模式检查金额和签名，确认正确后通过显式 \texttt{--broadcast} 参数提交交易：
\begin{lstlisting}[language={}]
python '.\split_test_coins .py' --broadcast 2>&1 |
    Tee-Object split_output.txt
\end{lstlisting}
\shot{split_terminal.png}{91mm}{分币脚本输出：十份各 14900 sat，脚本验证通过且 HTTP 200}
广播返回的分币交易哈希为：
\txhash{86c6408a8b95ca4bd7fa71c3f2ef362016d3fc4ef01e59fc62742d1c5eeae2d6}
随后通过交易浏览器检查输出列表和确认状态。交易最初显示 \texttt{Unconfirmed}，说明已传播但尚未入块；等待后获得确认，所在区块为 \textbf{5151331}。只有链上确认才说明本次分币交易已经被收录，\texttt{HTTP 200} 本身不能替代确认。

分币后，领取交易的原输出已被消费，新的十个输出编号为 0 至 9，均属于本人地址。后续 Ex1 必须引用这笔新交易的哈希，不能再次引用已花费的领取输出。

\newpage
\shot{split_confirmed.png}{188mm}{分币交易已确认：十个输出均为本人地址，每份 0.000149 BTC}
\shot{split_block.png}{34mm}{分币交易被收录至区块 5151331}

\newpage
\subsection{补全 P2PKH 锁定脚本与解锁脚本}
\texttt{ex1.py} 的两个核心 TODO 按课程要求在操作码级别完成，没有用 \texttt{to\_scriptPubKey()} 直接替代锁定脚本实现：
\begin{lstlisting}
def P2PKH_scriptPubKey(address):
    if not isinstance(address, P2PKHBitcoinAddress):
        raise ValueError('P2PKH address required')
    return [OP_DUP, OP_HASH160, bytes(address),
            OP_EQUALVERIFY, OP_CHECKSIG]

def P2PKH_scriptSig(txin, txout, txin_scriptPubKey):
    from config import my_private_key, my_public_key
    signature = create_OP_CHECKSIG_signature(
        txin, txout, txin_scriptPubKey, my_private_key)
    return [signature, my_public_key]
\end{lstlisting}
以上摘录仅将报错提示译为英文，便于排版，核心实现与实际脚本相同。\texttt{bytes(address)} 获取 P2PKH 地址解码后的 20 字节公钥哈希，而不是地址字符串的字符编码。

签名辅助函数对待发送交易计算签名摘要，并在签名后附加 \texttt{SIGHASH\_ALL} 字节。\texttt{scriptSig} 的两个元素依次为签名、公钥，对应上一节介绍的栈执行顺序。

发币输入选取分币交易的第 \textbf{0} 号输出，金额为 \textbf{14900 sat}。这里的 0 表示分币交易的首个输出，不表示“第一次交易”。主程序从 \texttt{transactions.py} 读取成功广播的分币哈希，并设置：
\begin{lstlisting}
amount_to_send = Decimal('0.000139')
utxo_index = 0
txid_to_spend = args.txid or split_txid
\end{lstlisting}
目标输出锁定到课程配置中的收款地址：
\txhash{mv4rnyY3Su5gjcDNzbMLKBQkBicCtHUtFB}
\[
14900 = 13900+1000\quad (\mathrm{sat}).
\]
代码在本地校验通过后才允许广播。该过程用本人的签名解锁旧输出，同时为收款方创建新的锁定输出；收款方后续花费这笔资金时需要其自己的私钥。

\newpage
\subsection{发币运行结果}
分币确认后，先运行 \texttt{python ex1.py} 检查预览，再实际广播并保存终端输出：
\begin{lstlisting}[language={}]
python '.\ex1.py' --broadcast 2>&1 |
    Tee-Object ex1_output.txt
\end{lstlisting}
\shot{send_terminal.png}{91mm}{Ex1 发币终端结果：输出 13900 sat，手续费 1000 sat，广播成功}
本次实际发币交易哈希为：
\txhash{c49a4a8c6772d1b5fc51814b20cbc5e57dce8d58c8896b00226a4e94c05dd4e6}
\par\noindent\begin{tabularx}{\linewidth}{@{}l X@{}}
\toprule 检查项 & 结果 \\
\midrule
交易输入 & 分币交易的第 0 号输出，14900 sat \\
收款输出 & 课程地址，13900 sat（0.000139 BTC）\\
手续费 & 1000 sat（0.00001 BTC）\\
实际大小与费率 & 192 bytes，约 5.21 sat/vB \\
本地验证 & \texttt{VerifyScript: PASS} \\
广播返回 & \texttt{HTTP: 200 OK}，返回哈希与本地交易哈希一致 \\
链上结果 & 已确认，区块 5151332 \\
\bottomrule
\end{tabularx}
发币输出仅有一份，输入与输出间的 1000 sat 差额全部作为手续费。其他九个分币输出没有被本次发币交易引用。

\newpage
\subsection{发币链上确认截图}
\shot{send_confirmed.png}{190mm}{发币交易已确认，输入 0.000149 BTC，向课程地址输出 0.000139 BTC}
截图显示发币交易获得确认。最终公开 API 核对结果为 \texttt{confirmed=true}，确认区块 \textbf{5151332}，与浏览器截图一致。网页中的输入解锁脚本及输出锁定脚本也对应本次 P2PKH 交易。

\newpage
\begingroup\small
\section{结果汇总与分析}
\par\noindent\begin{tabular*}{\linewidth}{@{\extracolsep{\fill}}l r r r r@{}}
\toprule 阶段 & 本人输入/sat & 输出总额/sat & 手续费/sat & 确认区块 \\
\midrule
领取 & --- & 150524 & --- & 5148688 \\
分币 & 150524 & 149000 & 1524 & 5151331 \\
发币 & 14900 & 13900 & 1000 & 5151332 \\
\bottomrule
\end{tabular*}
领取行只列出水龙头交易中属于本人的输出；水龙头自身的输入、其他输出及其手续费不计入本人的实验花费。分币和发币两笔交易的手续费合计为 $1524+1000=2524$ sat。本人保留九份分币输出，共 $9\times14900=134100$ sat，课程地址收到 13900 sat：
\[
150524=134100+13900+2524\quad (\mathrm{sat}).
\]
这说明资金流向与手续费满足守恒关系。三笔交易的完整哈希保存在 \texttt{transactions.py}，最终查询记录保存至 \texttt{chain\_status.txt}。
\subsection{遇到的问题与解决方法}
\begin{enumerate}
\item \textbf{DLL 加载失败。} 实际问题是加载了其他软件的不匹配 DLL；让当前进程优先找到课程目录中已验证的 64 位 DLL 后恢复正常，无需修改系统全局配置。
\item \textbf{旧查询接口找不到交易。} BlockCypher 未查到领取记录，改用助教提示的 testnet3 浏览器核对，并使用 Esplora 广播接口；不能因一个浏览器查不到就断定资金未到账。
\item \textbf{输出编号容易误填。} 依据交易输出列表确定领取输出为 1、分币后使用输出为 0。
\item \textbf{浮点金额不适合交易输出。} 使用 \texttt{Decimal} 计算并转为整数 sat，同时检查输出总额小于输入，留出手续费。
\item \textbf{广播与确认混淆。} 广播后先显示未确认，等待矿工打包后再保存确认截图；同一输出不能因等待而重复花费。
\end{enumerate}
\section{实验总结}
本次实验完成测试币领取记录核对、十输出分币、P2PKH 脚本编写、交易签名与广播，并确认两笔新交易已入块。通过实际操作，理解了 UTXO 的定位和整体花费规则、公钥哈希与签名的双重验证，以及手续费由输入输出差额形成的机制。

数字签名解决花费授权问题，节点对 UTXO 和交易规则的验证及区块确认共同维护资金记录的一致性。本地脚本验证、接口接受和链上确认是不同阶段。

\section*{参考资料}
\begin{enumerate}
\item 课程 \texttt{Ex1.pdf}
\item Blockstream Esplora API：\url{https://github.com/Blockstream/esplora/blob/master/API.md}。
\item 本次实验的 \texttt{transactions.py}、运行输出、截图与最终链上查询记录。
\end{enumerate}
\endgroup
\end{document}
